\documentclass[conference]{IEEEtran}

\usepackage{amsmath}
\usepackage{array}
\usepackage{booktabs}
\usepackage{microtype}
\usepackage{tabularx}
\usepackage{url}
\usepackage[hidelinks]{hyperref}

\newcommand{\code}[1]{\texttt{\small #1}}
\newcommand{\tbd}{\textbf{[\,--\,]}}  % result pending re-run
\newcolumntype{Y}{>{\raggedright\arraybackslash}X}
\renewcommand{\arraystretch}{1.05}

\title{Passive Pre-Association Wi-Fi Access-Point Selection}

\author{
\IEEEauthorblockN{Alistair}
\IEEEauthorblockA{Wi-Fi AP Selection Research Project}
}

% IEEEtran typesets the abstract in bold. Restore normal weight, keeping the
% italic "Abstract---" lead-in that the format expects.
\makeatletter
\let\@IEEEoldabstract\abstract
\def\abstract{\@IEEEoldabstract\normalfont\normalsize}
\makeatother

\begin{document}
\maketitle

\begin{abstract}
A Wi-Fi client in overlapping coverage must select an access point (AP)
before it can directly measure the service that each candidate would provide.
Signal strength describes link quality, but not contention or load, and may
therefore favor a nearby but congested AP. This project investigates whether a
client can make a better decision using only passive observations collected
before association, without probe requests or assistance from the network.
This report defines the experimental foundation of that investigation: a
controlled ns-3 simulator that creates multi-AP deployments, records decoded
frames and physical carrier-sense activity, and replays each fixed topology
once per candidate AP, changing only the AP the candidate joins. Every replay
of a topology therefore shares one identical pre-association observation while
yielding its own measured throughput. The simulator separates physical-topology
randomness from radio randomness and produces a structured choice-set dataset.
The resulting method supports a reproducible comparison between conventional
signal-based selection and selection based on richer passive evidence.
\end{abstract}

\section{Problem Formulation}
\label{sec:problem}

Throughput depends on channel contention, competing stations, transmission
rates, retries, and the airtime consumed by neighboring basic service sets.
Received signal strength reflects none of them directly. Those effects do,
however, leave traces a client can hear without transmitting: the frames
themselves, and the time the medium spends busy. This work asks whether such traces are enough:
\emph{can a client choose the AP that will provide the greatest throughput
using only what its radio can passively observe before association?}

Existing load-aware methods \cite{dely2014bestap,adame2021channel} make their
load estimate after association, from AP-provided reports, or from a network
controller. A client making its first choice has none of these: only beacons
and some data frames, in which AP distance and offered traffic appear only
through their effects.

Let $\mathcal{A}=\{a_1,\ldots,a_N\}$ be the APs in one deployment and let $O$
be the candidate station's passive observation before it joins any AP. If
$y_i$ is the throughput obtained after joining $a_i$, the desired decision is

\begin{equation}
    a^* = \underset{a_i\in\mathcal{A}}{\operatorname{arg\,max}}\; y_i,
\end{equation}

using a decision function whose inputs are derived only from $O$. In a live
network only one $y_i$ is ever measurable: the station joins a single AP and
observes that AP's throughput, while the values of $y_j$ for every AP it did
not join are counterfactual. The simulator recovers all $N$ values by replaying
the same scenario once per target AP, changing only the AP the candidate joins.
Every choice set therefore carries its true $a^*$.

The replay is sound only if $O$ is identical across all $N$ runs of a scenario.
If any part of the observation varied with the AP eventually joined, a decision
function could reach the right answer through a simulation artifact rather than
through the network. This is checked rather than assumed. The assumptions that
do remain are stated next.

\section{Experimental Scope and Assumptions}

Simulator 1 isolates the relationship among path quality, background load,
and passive evidence. It does not attempt to reproduce every feature of a
deployed WLAN. Table~\ref{tab:assumptions} states the principal assumptions so
that conclusions can be interpreted within the intended scope.

\begin{table}[!ht]
\caption{Controlled assumptions of Simulator 1.}
\label{tab:assumptions}
\centering
\footnotesize
\begin{tabularx}{\columnwidth}{@{}p{0.27\columnwidth}Y@{}}
\toprule
\textbf{Aspect} & \textbf{Current assumption} \\
\midrule
Deployment & Static, single-floor, no walls; one stationary candidate and
stationary background stations. \\
Radios & Homogeneous IEEE 802.11n devices at 5 GHz with 20 MHz channels and
Minstrel-HT rate control \cite{ieee80211}. \\
Propagation & Yans Wi-Fi channel; log-distance loss (exponent 3.0, 46.6777 dB
reference loss at 1 m) followed by Nakagami fading. \\
Traffic & Uplink UDP; constant background offered rates and a saturated
candidate. Wired backhaul is not a bottleneck. \\
Observation & Passive and client-side; no probe request, association, or
network-provided load report during feature collection. \\
Target & Mean candidate uplink throughput from successful association until
the simulation stops; failed association is assigned zero. \\
\bottomrule
\end{tabularx}
\end{table}

The environment is static and open, so it excludes mobility, walls, hardware
diversity, and application demand that varies over time. The label is uplink
UDP throughput, so it says little about TCP goodput, downlink performance, or
latency.

The choice set contains only APs discovered during the projected scan. In
passive mode, discovery requires at least one decoded beacon, and observations
with fewer than two discovered APs are excluded because no selection decision
exists.

\section{Simulator Design}

The simulator is built on ns-3.45 and its Wi-Fi models \cite{ns3models}. APs
sit 30 m apart in a fixed layout: two form a pair, three an equilateral
triangle, and four, six, and eight fill $2\times2$, $3\times2$, and $4\times2$
grids. Channels 36, 40, 44, and 48 are assigned in AP-index order and reused
once a deployment exceeds four APs. Because the layout is fixed as well, co-channel
separation takes a single value per deployment size --- 30 m at eight APs, 42 m
at six, and no reuse at or below four --- which ties channel reuse to AP count.
Section~\ref{sec:next} returns to this.

Each scenario draws 2, 3, 4, 6, or 8 APs and 2, 3, or 4 background stations
per AP. A quarter of scenarios spread those stations uniformly. The
rest place one or more hotspot APs, with fewer hotspots more likely and the
count capped at $\lceil N_{AP}/3\rceil$. Around those hotspots 70\% of
background stations fall in 10 m discs, and the remainder cover the
footprint. Mixing uniform and deliberately clustered regimes follows common
practice in WLAN studies \cite{wong2018daga,zhang2020placement}, and matters
here because under uniform load every AP is about equally busy and the choice
collapses back to signal strength. Each station
associates with its nearest AP, and every background station in a scenario
offers the same load, drawn uniformly at random from 1.0, 2.5, 4.0, and
6.0 Mbit/s.
Candidate positions are drawn to keep the decision non-trivial: 70\% near a
boundary between neighboring APs, 30\% at 4--16 m from a single AP. These are
sampling policies rather than radio parameters, and they are recorded with the
dataset because they set its difficulty.

Background traffic starts at $t=1$ s. Passive recording runs until $t=5.5$ s,
covering 4.5 s of loaded traffic; the candidate associates at $t=6$ s and the
run ends at $t=18$ s. The 0.5 s gap between recording and association is a
guard: ns-3 performs target-dependent polling just before a join, and without
the gap that polling would land inside the observation window. After joining,
the candidate offers 100 Mbit/s in 1250-byte packets, which keeps it backlogged
so that measured throughput reflects the service the network gives it rather
than the limits of its own traffic generator.

\subsection{Matched Counterfactual Replays}

For each sampled topology and radio seed, the experiment runs once per target
AP. AP and station positions, background traffic, the candidate's position, and
every pre-association device stay fixed; only the AP selected at $t=6$ s
changes. A \code{topologySeed} controls background-station placement and an
independent \code{rngSeed} controls fading, contention, and rate adaptation,
with five radio seeds generated per topology. Separating the two streams is
what allows stochastic radio behavior to be repeated without the physical
deployment shifting underneath it.

\subsection{Validation}

Validation runs before any model is fitted, so that a good accuracy number
cannot mask a defective data-generating process.

Three properties carry the experimental design, and each is checked
automatically. Repeating a seed reproduces throughput exactly, so any run can
be regenerated. Changing only the radio seed varies the outcome while every
physical position stays fixed. And the recorded frame and channel-busy
traces are byte-identical across all target-AP replays of a matched scenario
--- the check that enforces the constant $O$ required in
Section~\ref{sec:problem}, and the most important of the three.

These establish internal consistency, not that the simulator behaves like a
real network.

\section{Observation and Dataset Construction}
\label{sec:obs}

\subsection{Recording and Projection}

The candidate carries one receive-only scanner PHY per occupied channel, plus
a separate PHY reserved for association. The scanner PHYs never transmit and
never associate; the association PHY is parked on unused channel 149 until the
decision time, so the AP the candidate will eventually join cannot alter what
is heard, or consume random draws, while the observation is being collected.
At $t=5.5$ s the scanner PHYs move to the parking channel as well.

Watching every channel at once is not a claim that real clients have several
radios. It is a recording device: it captures a superset of what any single
radio could have heard, from which a realistic single-radio observation is cut
afterwards.

Two traces are recorded. The monitor-sniffer trace holds every frame a scanner
decodes, with time, frequency, BSSID, transmitter, frame category, beacon and
retry flags, length, RSSI, noise, airtime, and PHY rate. The PHY-state trace
gives the fraction of each 102.4 ms beacon interval in which carrier sense
marked the channel busy. This clear-channel-assessment (CCA) measurement
registers energy a frame-only view misses: interference, and transmissions too
distant to decode. Real chipsets expose such a counter, though whether a driver
surfaces it varies by platform.

The superset is then projected onto one passive radio. That radio rotates
through the four channels in random order, dwelling 110 ms on each with
0.15 ms to retune, and repeats the cycle for as much of the window as it
fits, so every channel is revisited many times before the decision. Frames
and CCA intervals falling outside a dwell are discarded. A 110 ms dwell
spans just over one 102.4~ms beacon interval, so each visit yields about one
beacon from every AP on that channel. RSSI is then perturbed by uniform
error within $\pm3$ dB and quantized to 1 dB. Both steps run after ns-3 from
a recorded scan seed, so scan policy and measurement noise can be varied
without repeating a simulation.

\subsection{From Trace to Training Data}

The projected trace is a long list, one entry per decoded frame, with the
channel-busy series recorded alongside it. Turning that list into training
data means collapsing it once per discovered AP. For a given AP the builder
first selects the parts of the trace that bear on it --- frames sent by or to
its BSSID, and activity on the channel it occupies --- which yields a shorter
list, not yet a single description. It then reduces that shorter list to a
fixed-size summary of the AP, and it is the summary, not the frames, that a
model ever sees. Four discovered APs therefore give four summaries over
identical evidence, differing only in which AP each describes and in the
throughput that AP delivered when it was the target of a replay.

That reduction is carried out twice, in two different ways. The first
produces a fixed vector of summary statistics --- these are the
\code{feat\_} columns, and they are hand-chosen. The second keeps time: the
window is cut into equal bins and each measurement is recorded separately in
every bin, producing a sequence per option instead of a vector. Tabular
models are trained on the first, the transformer on the second, and
Section~V compares them.

The summary vector carries three kinds of feature. AP-level features describe
that BSS itself: RSSI, beacons, observed clients, frames, airtime, retries,
and PHY rates. Channel-level features describe the option's channel as a
whole, decoded traffic and CCA busy time together. Relative features place the
option against its rivals: how many options there are, its RSSI rank, and its
margin over the best alternative.

Ground truth is kept but quarantined. True distance, configured load, station
count, positions, and seeds are all recorded, because they are needed to check
that the observable features behave and to break results down by scenario type.
None may reach a model. The two are separated by name: only \code{feat\_}
columns are admissible inputs, everything with a \code{gt\_} prefix is not,
and the loader enforces the split rather than trusting discipline.

Rows are grouped twice over. Each physical topology carries a
\code{topology\_id}, and each of its five radio-seed realizations carries its
own \code{group\_id}. Rows sharing a \code{group\_id} were produced from one
identical observation and differ only in which AP was joined. Partitions must
therefore be cut by \code{topology\_id}: the five realizations of a layout are
near-duplicates of each other, and letting them fall on both sides of a
train/test split would inflate every result.

\section{Models and Evaluation Protocol}

Two observations are compared. The \emph{rotating} observation is the
single-radio projection of Section~\ref{sec:obs}, in which one radio cycles
through the channels for the whole window. The \emph{continuous} observation
keeps the parallel-scanner superset instead, so every channel is watched at
once. One radio cannot do this, so the continuous condition is an upper bound
on what observation could provide rather than a deployable alternative. Both
carry the same RSSI realism, and both are binned into per-option sequences.

The dataset holds \tbd{} AP options across \tbd{} observations of \tbd{}
topologies, a median of \tbd{} discovered APs per decision. Choosing the best
AP rather than at random is worth \tbd{} against \tbd{} Mbit/s on average, and
the mean spread between the best and worst option in a choice set is \tbd{}
Mbit/s.

Four families of decision function are compared. Heuristics score an option
directly: strongest RSSI, least busy channel, and RSSI minus channel-busy
fraction. Tabular models take one feature vector per option, under three
encodings of the same sequence --- its window means, a set of summary
statistics per measurement (mean, standard deviation, minimum, maximum, last
value, and least-squares slope), and the whole sequence flattened with bin
order preserved. A ridge regression takes the flattened sequence as a linear
control. The transformer encodes each AP's sequence with a learned positional
embedding, then compares APs through a permutation-invariant set stage, and is
trained twice: once to predict throughput, once to rank options directly.

Every model sees the same five splits, cut by \code{topology\_id} and
stratified by AP count. Metrics are top-1 accuracy with a 0.5 Mbit/s tie
tolerance, mean regret in Mbit/s against the oracle choice, and $R^2$ where a
model predicts throughput. Significance is a paired $t$-test across the five
splits.

\section{Results}

Table~\ref{tab:results} gives top-1 accuracy for every model under both
observations. Heuristic reference points are identical in the two columns,
since they read the same features: random selection reaches \tbd{}, least busy
channel \tbd{}, strongest RSSI \tbd{}, and RSSI minus busy \tbd{}.

\begin{table}[!ht]
\caption{Top-1 accuracy by model and observation. Within a column every model
receives identical input.}
\label{tab:results}
\centering
\footnotesize
\begin{tabularx}{\columnwidth}{@{}Yrrr@{}}
\toprule
\textbf{Model} & \textbf{Rotating} & \textbf{Continuous} & \textbf{Gain} \\
\midrule
Ridge, flattened sequence   & \tbd & \tbd & \tbd \\
MLP, flattened sequence     & \tbd & \tbd & \tbd \\
MLP, summary statistics     & \tbd & \tbd & \tbd \\
GBM, summary statistics     & \tbd & \tbd & \tbd \\
GBM, flattened sequence     & \tbd & \tbd & \tbd \\
Transformer                 & \tbd & \tbd & \tbd \\
\bottomrule
\end{tabularx}
\end{table}

Three comparisons are read off this table. The first is which model wins under
each observation, and by how much, given that every model in a column receives
identical input. The second is what widening the observation is worth, model
by model, which the gain column reports. The third is whether the winner
changes between the two columns.

Two ablations accompany it. Giving every tabular model the mean of the other
options and its own deviation from that mean supplies the cross-option
comparison the transformer's set stage performs, and isolates how much of any
transformer margin is attributable to it. Shuffling the time bins within each
option, leaving bin position intact, removes temporal order while preserving
everything else, and isolates aggregation from tracking.

\section{Does the Sequence Carry Anything?}

The matched design allows the label's variance to be split directly. Across
cells of one topology paired with one target AP, each realized under five
independent radio seeds, the intraclass correlation separates what is fixed
when the topology is drawn from what varies with the radio alone. A high value
bounds what any time-sensitive model can contribute, because it says the
outcome is settled before the observation begins; a low one leaves room for
the trajectory to matter. The measured value is \tbd{}, against a between-cell
variance of \tbd{} and a within-cell variance of \tbd{}.

The shuffle ablation is the independent check on the same question. If order
carries information, destroying it must cost accuracy; if the sequence is
repeated sampling of a stationary environment, it will not. The two
instruments are reported together because agreement between them is what makes
either conclusive.

\section{What Simulator 2 Must Change}
\label{sec:next}

The requirement is measurable rather than a matter of judgement.
Recomputing the intraclass correlation on a successor dataset is a direct test
of whether the question can be asked at all: unless it falls well below 0.9,
no amount of model engineering will make a temporal model useful, because
there will again be nothing for it to track. What must change is therefore not
the model but the environment --- load that arrives and departs, interference
that appears part-way through the window, stations that join and leave ---
and it must be visible in the passive trace while remaining invisible to a
window average.

Four smaller decisions follow from the present work. Channel assignment by AP
index, over a fixed layout, leaves co-channel separation with no variance and
ties reuse to deployment size; sampling the assignment instead would decouple
the two. AP spacing of 30 m pairs an obstructed-indoor path-loss exponent with
an open-field geometry, and the two should be reconciled. The BSS Load element, which a passive client can
read from a single beacon, is absent here and is the strongest baseline this
comparison currently lacks. And the restriction to complete silence, rather
than to the pre-association period, is a scope choice that could be relaxed
to admit probe requests.

One practical consequence deserves separating from the rest. Since observation
budget moved accuracy further than any change of model, and one radio cannot
watch every channel at once, the lever a real client actually holds is how
long it may scan before deciding. Dwell time and number of passes are
therefore worth treating as experimental variables in their own right, rather
than as fixed properties of the projection.

\begin{thebibliography}{00}

\bibitem{dely2014bestap}
P. Dely, A. Kassler, L. Chow, N. Bambos, N. Bayer, H. Einsiedler, and
C. Peylo, ``BEST-AP: Non-intrusive estimation of available bandwidth and its
application for dynamic access point selection,'' \emph{Computer
Communications}, vol. 39, pp. 78--91, 2014,
doi: 10.1016/j.comcom.2013.11.006.

\bibitem{adame2021channel}
T. Adame, M. Carrascosa, B. Bellalta, I. Pretel, and I. Etxebarria, ``Channel
load aware AP/Extender selection in home WiFi networks using IEEE
802.11k/v,'' \emph{IEEE Access}, vol. 9, pp. 30095--30112, 2021,
doi: 10.1109/ACCESS.2021.3059473.

\bibitem{wong2018daga}
W. Wong and S.-H. G. Chan, ``Distributed joint AP grouping and user
association for MU-MIMO networks,'' in \emph{Proc. IEEE INFOCOM}, 2018,
pp. 252--260, doi: 10.1109/INFOCOM.2018.8486284.

\bibitem{zhang2020placement}
Y. Zhang and L. Dai, ``Joint optimization of placement and coverage of access
points for IEEE 802.11 networks,'' in \emph{Proc. IEEE ICC}, 2020,
pp. 1--7, doi: 10.1109/ICC40277.2020.9149098.

\bibitem{ieee80211}
IEEE Standard for Information Technology---Telecommunications and Information
Exchange between Systems---Local and Metropolitan Area Networks---Specific
Requirements---Part 11: Wireless LAN Medium Access Control (MAC) and Physical
Layer (PHY) Specifications, IEEE Std. 802.11-2024, 2025.

\bibitem{ns3models}
ns-3 Project, \emph{ns-3 Model Library, Release 3.45}, 2025. [Online].
Available: \url{https://www.nsnam.org/docs/release/3.45/models/}

\end{thebibliography}

\end{document}
